\section{What remains}\label{sec:open}

Two questions remain open after this memo: one about the field, and the parent paper's Problem~9.1(c) in a sharper form.

\begin{problem}[Characteristic three]\label{prob:char3}
Over a field of characteristic~$3$, do the classes $\{0,1,3\}$ and $\{2,4,5\}$ of $\Z/6$ carry a code of a
class-uniform rule of dimension~$25$?
\end{problem}

Both rational certificates of \cref{thm:fields} have scales divisible by $3$ ($N=1008$ and $P=3$; $N=187200$ and
$P=15$), so neither certificate reduces modulo~$3$. A positive answer would make the gain of \cref{prop:F2} independent
of the field; a negative one would make it depend on the characteristic.

\begin{problem}[Problem~9.1(c) of the parent paper, for a bounded number of classes]\label{prob:bounded}
Fix $M$, for instance $M=2$. Is the excess $D(\mathcal P)-(d^2+Md)/2$ bounded over all $d$ and all partitions
$\mathcal P$ of $\Z/d$ into $M$ classes? If not, find a family of such partitions and codes, with a description that
can be analysed, in which it grows.
\end{problem}

By \cref{cor:excess}, blow-ups trade a growing number of classes for a growing excess at a fixed ratio, so they do not
bear on \cref{prob:bounded}. A growing excess at a fixed number of classes would be a coding gain that grows with $d$ and
is native to masks without an almost-period, the first of its kind in the one-hop model; it would still be consistent
with a positive answer to Problem~9.1(a), which concerns gains of order~$d^2$.

The steps from the one-hop model to machines \cite[Section~7.4]{SeventeenthPaper2026} change as follows. Depth
reduction, the second step, is priced in the one-hop model over finite fields by \cref{thm:any-advice}; for machines it
still needs what \cref{rem:depth} lists. The first step, a version of \cref{thm:TP} for the cone systems of real
schedules, and the third, many hops, are unchanged. Statement~(T) remains open.
